\section{Results}
\label{sec:results}

\begin{figure*}[t]
\centering
\includegraphics[width=\linewidth]{figures/A11_gaia_harness_interaction.pdf}
\caption{\textbf{On the one task where the harness really varies, the best harness depends on the model.} GAIA is the only task whose optimizer harness has more than two levels: every model ran under \oc, \goose\ and \miniswe\ as well as its own native harness. Its measured-zero baseline makes gain here the raw held-out score. \textbf{(a)} One line per model over a fixed harness order; the lines cross, so no harness holds its rank across models. \textbf{(b)} Each model's native harness against its own best shared one, against the task's resolution band. Both GPT models are far better under \codex\ ($+0.179$ and $+0.131$); the two Claudes and Kimi sit within a band or two of zero either way.}
\label{fig:gaia-harness}
\end{figure*}

Table~\ref{tab:gain-by-task} compares the core experimental configurations
across tasks on normalized gain. Appendix~\ref{app:master} reports the same
results in Table~\ref{tab:master}, together with run ranges and properties of
the generated harnesses by task.


\subsection{Distinguishing frontier models}
\label{sec:rq1}

Model choice has a larger effect than coding-harness choice. Holding task and
harness fixed, changing the optimizer model moves gain by \ModelSwap\ on
average; holding task and model fixed, changing the harness moves it by
\HarnessSwap. The model contrast is therefore about
$\AxisRatio\times$ larger. Both exceed the task resolution bands, although the
harness contrast does so narrowly.

The extremes separate more clearly than the middle. The strongest configuration
captures roughly two thirds of the available OfficeQA headroom and half of the
BrowseComp-Plus headroom, while the weakest is unresolved from zero on
BrowseComp-Plus and Terminal-Bench. Differences among intermediate
configurations are often smaller than their round-to-round variation, supporting
tiers rather than a complete ranking. Figure~\ref{fig:showdown} shows the run-level spread and the corresponding
task-adjusted model effects.

\subsection{Tracking model progress}
\label{sec:ladder}

We test sensitivity to model progress using two OfficeQA release series, holding
the target model, seed, budget, and coding harness fixed within each series.
Across \NumLadderRungs\ GPT releases, gain rises monotonically from
\LadderGptLo\ to \LadderGptHi, with three of its four steps exceeding the task's
resolution band. Across \NumLadderRungs\ Claude Opus releases, gain ranges from
\LadderClaudeLo\ to \LadderClaudeHi; gain is non-monotonic, but the first-to-last spread exceeds the task's resolution band.


\subsection{Where current optimizers fall short}
\label{sec:process}

Trajectory-derived measures characterize how optimizers search; they are not
independent measures of held-out capability.

\paragraph{Broader search is associated with greater gain.}
We identified eight harness levers on OfficeQA before examining the other tasks:
prompt, context management, step cap, retry and timeout policy, tool schema,
answer extraction, retrieval policy, and reasoning effort. The fraction touched
during search is positively associated with gain on every task, with
$\rho$ ranging from $\LeverRhoLo$ to $\LeverRhoHi$
(Figure~\ref{fig:levers}). No other process measure we investigated had the same direction on
all four tasks.

This is \emph{exploration} rather than final candidate breadth. An optimizer may inspect or modify
a lever without retaining the change: one configuration touched three quarters
of the levers and made seven edits but shipped the original seed. Breadth is also
correlated with total modification volume, so the data do not isolate breadth
from search effort.


\paragraph{Trace reading is not associated with higher gain.}
The share of actions spent reading evaluation output is negatively associated
with gain across the four tasks, from $\TraceRhoHi$ to $\TraceRhoLo$.
Optimizers rely mainly on per-case score summaries: detailed trace spans were
requested only \NSpanCalls\ times by \NSpanCells\ of the
\NumScoredCells\ cells. This does not establish that diagnosis is unnecessary.
For these tasks, per-case summaries may localize failures well enough that reading
full traces does not justify its context cost.

\paragraph{Case passes, not evaluation calls, bind.}
The development and validation partitions each permit \EvalCallCap\ evaluation
invocations and four full case passes. The median optimizer uses \EvalCallsUsed\ calls
(\EvalCallFrac) but \CasePassFrac\ of its case allowance, and
\NBudgetExhausted\ of \NBudgetCells\ cells exhaust at least one partition's case
budget. Thus the case allowance constrains search, whereas the call cap does not.

\paragraph{Visible validation scores are optimistic.}
Most cells in Figure~\ref{fig:overfit} fall below the identity line: the
submitted candidate's test score is lower than the best validation score observed
during search. The held-out partition is therefore necessary to measure realized
gain. The figure cannot distinguish selection-induced overfitting from a
validation--test mismatch, so we claim only that the visible best score is
optimistic.

\subsection{The effect of the optimizer's own harness}
\label{sec:rq3}

Across the \NHarnessPairs\ model--task pairs evaluated under both conditions,
the shared harness wins \CommonWins, the native harness wins \NativeWins, and
\HarnessTies\ are tied. Native harnesses therefore have no consistent advantage.
Restricting evaluation to a model's native tooling would not provide a reliable
estimate of its harness-optimization ability.

In \NHarnessPairsResolved\ pairs, the difference exceeds the task's resolution
band, but the direction varies by model and task. The aggregate near-tie
therefore reflects heterogeneous effects rather than uniformly small ones.

Figure~\ref{fig:gaia-harness} shows where the effect does live, on GAIA --- the
only task whose harness axis has more than two levels. The harnesses do not hold
their rank across models, and the native-harness advantage is concentrated rather
than absent: both GPT models are four to five resolution bands better under
\codex, while the two Claudes and Kimi sit within a band or two of zero. Reading
only the shared harnesses would miss this, and would suggest that sensitivity to
the harness scales with model capability; adding each model's native harness
removes that pattern.
